\section{Uniform expected-depth hardness}
\label{sec:uniform-hardness}

The weighted reduction of Section~\ref{sec:hardness} also yields
uniform-prior hardness, but a direct replacement of a weight by
ordinary points requires a normalization argument. We prove an exact
replacement theorem. It preserves the weighted optimum up to a
computable additive term, despite the additional private queries.
Throughout this section, raw uniform cost means the sum of the leaf
depths, before division by the number of hypotheses.

\subsection{A rotation for a private singleton}

\begin{lemma}[Uniform singleton rotation]\label{lem:uniform-singleton-rotation}
Suppose a proper identifying tree on a uniformly weighted state $S$
first asks a supplied query isolating $x\in S$, and the other child
first asks a query $r$ that is proper on $S\setminus\{x\}$.
These two queries can be exchanged, placing the singleton query
in the answer branch of $r$ containing $x$, without increasing
the raw cost.
\end{lemma}
\begin{proof}
Let $A,B$ be the two nonempty answer classes of $r$ on
$S\setminus\{x\}$, and choose their names so that $x$ has the
answer of $A$. In the original tree, $x$ has relative depth one
and every hypothesis in $A\cup B$ passes through both queries.
Ask $r$ first, put the private singleton query in its
$A\cup\{x\}$ branch, and attach the original subtrees on $A$ and
$B$. The depth of $x$ increases by one, the depth of every point
of $B$ decreases by one, and all other depths are unchanged.
The change in raw cost is $1-|B|\le0$. The resulting queries are
proper and are members of the original supplied family.
\end{proof}

If the two queries isolate different singleton points, the same
rotation has zero cost. The lemma concerns equal weights; it
does not assert the corresponding inequality for unequal priors.

\subsection{An exact geometric cloud replacement}

Consider the weighted singleton/pair linear forest instance of
Lemma~\ref{lem:pair-invoice}, with vertices $0,1,\ldots,N$ and
positive integer weights $w_0,\ldots,w_N$. Vertex $0$ is
distinguished. Order the nonzero vertices along each path and
component by component, so every edge joins consecutive indices. Define
\[
 n=\sum_{i=0}^N w_i,\qquad Q=4\max_i w_i,\qquad
 C_i=\{(Qi+j,(Qi+j)^2):0\le j<w_i\}.
\]
The point with $j=0$ is the anchor of cloud $C_i$.
For every nonzero vertex supply its whole-cloud singleton query;
for every linear forest edge supply its whole-cloud pair query.
For consecutive cloud indices $a,\ldots,b$, these are the negative
answers of
\begin{equation}\label{eq:uniform-cloud-membership}
 4y-4Q(a+b)x+Q^2(2a-1)(2b+1)=0.
\end{equation}
The value on the parabola factors as
$4(x-Q(a-\tfrac12))(x-Q(b+\tfrac12))$.
Every cloud has width less than $Q/4$, so the negative answer
selects exactly those clouds and excludes cloud $0$.

For every nonanchor point of abscissa $t$, including those of
cloud $0$, supply the private line
\[
 2y-4tx+2t^2-1=0.
\]
Its value on $(x,x^2)$ is $2(x-t)^2-1$; its negative answer
isolates precisely that point. No private line is supplied for
an anchor. All answers are strict and binary.
The whole-cloud singleton queries separate the anchors, and
the private lines distinguish every nonanchor, so the family
identifies every hypothesis. All hypotheses have uniform prior $1/n$.

Put
\begin{equation}\label{eq:uniform-cloud-offset}
 D(w)=\frac{(w-1)(w+2)}2,\qquad
 \Delta=\sum_{i=0}^N D(w_i).
\end{equation}
Algorithm~\ref{alg:uniform-cloud-lift} gives the construction and
the threshold conversion.

\begin{algorithm}[H]
\caption{Uniform cloud replacement of a weighted linear forest instance}
\label{alg:uniform-cloud-lift}
\KwIn{Vertices $0,\ldots,N$ with each path in order, component by component,
positive integer weights $w_i$, and a raw weighted threshold $k$}
\KwOut{Supplied-line instance with uniform prior and threshold
$(k+\Delta)/n$}
$Q\gets4\max_i w_i$; $n\gets\sum_iw_i$;
$\HH\gets\varnothing$; $\LL\gets\varnothing$\;
\For{$i=0,\ldots,N$}{
  \For{$j=0,\ldots,w_i-1$}{
    $t\gets Qi+j$; add $(t,t^2)$ to $\HH$\;
    \If{$j>0$}{
      Add $2y-4tx+2t^2-1=0$ to $\LL$\;
    }
  }
}
For each $i>0$, add \eqref{eq:uniform-cloud-membership}
with $(a,b)=(i,i)$ to $\LL$\;
For each linear forest edge $\{i,i+1\}$, add
\eqref{eq:uniform-cloud-membership} with $(a,b)=(i,i+1)$ to $\LL$\;
$\Delta\gets\sum_{i=0}^N(w_i-1)(w_i+2)/2$\;
\KwRet $(\HH,\LL)$, prior $1/n$, and threshold $(k+\Delta)/n$\;
\end{algorithm}

\begin{lemma}[Local uniform cloud costs]\label{lem:uniform-cloud-local}
The optimal raw local cost on one intact cloud of size $w$ is
$D(w)$. On two intact clouds of sizes $u,v$, it is
\[
 D(u)+D(v)+u+v.
\]
\end{lemma}
\begin{proof}
Within one cloud, every whole-cloud query is constant, and every
nonconstant query isolates one of its nonanchors. The only possible
proper code is a chain, whose sum of depths is
$1+2+\cdots+(w-1)+(w-1)=D(w)$; for $w=1$ the sum is zero.

On two clouds a nonconstant whole-cloud query separates the
clouds, attaining the displayed formula. Every other nonconstant
query isolates one nonanchor. Induct on $u+v$.
If a private point of the $u$ cloud is queried first, the
remaining two-cloud state has sizes $u-1,v$ and still has its
anchors and whole-cloud separator. Using induction, this choice
costs
\[
 u+v+D(u-1)+D(v)+(u-1)+v.
\]
Subtracting $D(u)+D(v)+u+v$ gives $v-1\ge0$, because
$D(u)-D(u-1)=u$.
The analogous difference for a private point of the $v$ cloud
is $u-1\ge0$. The base case $u=v=1$ has raw cost two.
When one size is one it has no private point, and the same
induction applies to the other cloud. Queries selecting other
clouds are constant or separate these two whole clouds, so the
argument covers the entire supplied family.
\end{proof}

\begin{theorem}[Exact uniform cloud transfer]\label{thm:uniform-cloud-transfer}
Let $\OPT_{\mathrm{raw}}^{\mathrm{weighted}}$ be the minimum
raw weighted cost of the singleton/pair linear forest instance.
For its uniform cloud replacement,
\begin{equation}\label{eq:uniform-cloud-transfer}
 n\OPT_{\mathrm E}^{\mathrm{cloud}}
   =\OPT_{\mathrm{raw}}^{\mathrm{weighted}}+\Delta.
\end{equation}
Equivalently, if $J(\mathcal M)$ has the pair and singleton masses
from Lemma~\ref{lem:pair-invoice}, the raw uniform optimum is
\[
 \Delta+\min_{\mathcal M}
 \left\{w_0(N-|\mathcal M|)
   +\mathcal C(J(\mathcal M))
   +\sum_{\{u,v\}\in\mathcal M}(w_u+w_v)\right\}.
\]
\end{theorem}
\begin{proof}
Start with a proper optimal uniform tree and follow the anchor
of cloud $0$. Suppose a private query for a nonanchor $x$ of a
nonzero cloud occurs on this path. Its own cloud anchor is still
on the path, since no private query isolates an anchor.
Repeatedly apply Lemma~\ref{lem:uniform-singleton-rotation} to
move this private query past the next query on the anchor-$0$
path. The next query is proper on the remaining state; constant
queries have already been removed. Private queries commute
at zero cost. Eventually a whole-cloud query separates $x$'s
anchor from the anchor of cloud $0$: otherwise the tree would
not identify both anchors. The point $x$ has the answer of its
own anchor to that query. After the rotation its private query
is in the hit branch, off the anchor-$0$ path.

Each such operation removes one external private query from
that path, introduces none, and moves it through only finitely
many existing path queries. Repetition therefore terminates.
Next commute each cloud-$0$ private query past all remaining
whole-cloud queries. Its point always has the answer of the
anchor of cloud $0$, so it stays on that path but moves after
the last whole-cloud query. This second process also terminates,
by decreasing the number of private/whole-cloud inversions.
Every rotation is cost nonincreasing.

Consequently all newly hit clouds are intact at their first
whole-cloud query: no private query for any of their points
has occurred on the path earlier. First-hit groups are whole
singleton clouds or whole adjacent pairs, and the pairs form
a matching $\mathcal M$. A query first hitting just one cloud
can be replaced by its supplied whole-cloud singleton query,
which has the same split on the current state. There are
$N-|\mathcal M|$ whole-cloud spine queries. Cloud $0$ contributes
$w_0(N-|\mathcal M|)$ to the spine cost and $D(w_0)$ afterwards.
Every other cloud contributes $D(w_i)$ by
Lemma~\ref{lem:uniform-cloud-local}; every pair also contributes
$w_u+w_v$. The first-hit positions charge their cloud masses.
Sorting these disjoint group masses nonincreasingly minimizes
that contribution, giving $\mathcal C(J(\mathcal M))$.
This is the asserted lower bound.

Conversely, for any matching, query its disjoint pairs and
unmatched singleton clouds in decreasing total cloud size.
Attach the optimal local codes of
Lemma~\ref{lem:uniform-cloud-local} and identify cloud $0$ by
its private chain at the end. This attains the matching invoice
plus $\Delta$. Taking the minimum proves equality.
\end{proof}

\subsection{Uniform-prior strong NP-completeness}

\begin{theorem}\label{thm:strong-uniform-average}
Deciding whether the optimal expected depth is at most a rational
threshold is strongly NP-complete under a uniform prior, even for
integer-coordinate points on a parabola, polynomially bounded
integer line coefficients, and $O(n)$ supplied strict binary lines.
\end{theorem}
\begin{proof}
Apply Algorithm~\ref{alg:uniform-cloud-lift} to the hard weighted
linear forest instances of Theorem~\ref{thm:strong-average}.
If the original number of hypotheses is $n_0=N+1$, that
construction has nonzero weights $O(n_0^3)$ and distinguished
weight $O(n_0^5)$. Thus the replicated point count
$n=\sum_iw_i=O(n_0^5)$ is polynomial in the restricted source size.
For a raw weighted threshold $k$, use the expected threshold
$(k+\Delta)/n$. Theorem~\ref{thm:uniform-cloud-transfer} gives
the yes/no equivalence for arbitrary adaptive trees.

The construction has $n-(N+1)$ private lines and at most
$2N$ whole-cloud lines, hence $O(n)$ supplied queries.
Since $Q\le4n$ and $N\le n$, its abscissae are $O(n^2)$,
and its ordinates and line coefficients are $O(n^4)$.
All are integers, and all answers are strict.
Both $k$ and $\Delta$ are bounded by fixed polynomials in $n$,
so the rational threshold has polynomially bounded numerator
and denominator. A proper identifying tree has $n-1$ internal
nodes and is a polynomial certificate; its raw uniform cost
is evaluated by summing leaf depths. This proves strong
NP-completeness.

\end{proof}

\begin{remark}[Limits of the reduction]
The cloud identity preserves exact hardness but adds the same
nonnegative internal cost $\Delta$ to every matching. An EPTAS for
the compressed weighted path family, such as
Theorem~\ref{thm:pair-eptas}, remains an EPTAS for these explicit
cloud instances: its additive error is at most
$\varepsilon\OPT_{\rm raw}^{\rm weighted}\le
\varepsilon(\OPT_{\rm raw}^{\rm weighted}+\Delta)$.
Thus this reduction does not establish constant-factor
inapproximability. That question requires a different construction.
\end{remark}
